\chapter{Business Cycles and Shocks}\label{ch:b2:business-cycles}

Output does not grow smoothly: it rises faster than its trend for a while, then slower, and sometimes falls. This chapter
measures these fluctuations, separates them from the trend, and builds the simplest equilibrium model in which they arise
from shocks to technology. The chart shows how volatile growth has been in rich countries over a century and a half.

\begin{openingchart}{Standard deviation of the annual growth of real GDP per person over the ten years ending in each year,
  1880--2020, per cent, logarithmic scale: United States, United Kingdom, Germany, and the median of seventeen rich
  countries.}{Jord\`a--Schularick--Taylor Macrohistory Database (rgdpmad); our computation}
\begin{tikzpicture}
\begin{axis}[omrecord, width=0.72\linewidth, height=0.42\linewidth, xmin=1880, xmax=2020, ymode=log, ymin=0.4,
    ymax=40, ytick={0.5,1,2,5,10,20}, yticklabels={0.5,1,2,5,10,20}, xticklabel style={/pgf/number format/1000 sep=},
    ylabel={Per cent}, legend style={at={(0.5,-0.18)}, anchor=north, legend columns=4, font=\scriptsize}]
  \addplot[omDef, thick] table[x=year, y=USA] {figdata/bachelor-2/28-business-cycles/out/vol.csv};
  \addlegendentry{United States}
  \addplot[omThm, thick, dashed] table[x=year, y=GBR] {figdata/bachelor-2/28-business-cycles/out/vol.csv};
  \addlegendentry{United Kingdom}
  \addplot[omMeth, thick, dotted] table[x=year, y=DEU] {figdata/bachelor-2/28-business-cycles/out/vol.csv};
  \addlegendentry{Germany}
  \addplot[black, very thick] table[x=year, y=median] {figdata/bachelor-2/28-business-cycles/out/vol.csv};
  \addlegendentry{median of 17}
\end{axis}
\end{tikzpicture}
\end{openingchart}

% figdata: vol_summary (median: 1913 2.34, 1950 6.20, 1975 2.21, 2000 1.54, 2007 1.13, 2019 0.89 (minimum), 2020 2.10; DEU 1950 27.29; USA 2007 1.11)
In the median country, growth of income per person had a standard deviation of 2.34 per cent over the ten years to 1913 and
6.20 per cent over the ten years to 1950, a decade of war and reconstruction; Germany's reached 27.29. It was 2.21 in the decade to 1975, 1.54 in
the decade to 2000 and 1.13 in the decade to 2007. The crisis of 2008--09 raised it; the calm decade 2010--19 brought it to 0.89, the
lowest of the whole period; and 2020 alone took it back to 2.10.
Fluctuations are a permanent feature, but their size has changed a great deal.

\begin{recall}
From the school volumes: the business cycle, recessions and expansions. From the Year 1 volume: aggregate demand and supply
shocks. From \cref{ch:b2:ramsey-model}: the planner's problem and the \omterm{def:b2:intertemporal-choice:euler-equation}{Euler equation}. From
\cref{ch:b2:three-equation-model}: demand shocks and monetary policy.
\end{recall}

\section{The facts}\label{sec:b2:business-cycles:facts}

\begin{definition}[Cyclical component]\label{def:b2:business-cycles:cyclical-component}
The \emph{cyclical component}\index{cyclical component} of a time series is its deviation from a smooth trend, usually of
its logarithm, so that it measures the per cent distance from trend.
\end{definition}

\begin{definition}[Procyclical variable]\label{def:b2:business-cycles:procyclical}
A variable is \emph{procyclical}\index{procyclical variable} if its \omterm{def:b2:business-cycles:cyclical-component}{cyclical component} is positively correlated with that
of output, \emph{countercyclical}\index{countercyclical variable} if negatively, and acyclical if the correlation is near
zero.
\end{definition}

The table summarises the American cycle from 1959 to 2009: the standard deviation of each \omterm{def:b2:business-cycles:cyclical-component}{cyclical component}, its size relative to
output's, its correlation with output's and its correlation with its own value one quarter earlier.

\begin{center}
\pgfplotstabletypeset[col sep=comma,
  columns={var,sd,rel,corr,ac1},
  columns/var/.style={string type, column name={Variable}, column type={l}},
  columns/sd/.style={column name={Std.\ dev.}, column type={r}, fixed, fixed zerofill, precision=2},
  columns/rel/.style={column name={Relative to output}, column type={r}, fixed, fixed zerofill, precision=2},
  columns/corr/.style={column name={Corr.\ with output}, column type={r}, fixed, fixed zerofill, precision=2},
  columns/ac1/.style={column name={Autocorr.}, column type={r}, fixed, fixed zerofill, precision=2},
  every head row/.style={before row=\toprule, after row=\midrule}, every last row/.style={after row=\bottomrule}]
  {figdata/bachelor-2/28-business-cycles/out/moments.csv}

{\small United States, HP \omterm{def:b2:business-cycles:cyclical-component}{cyclical components} (smoothing 1600) of quarterly data, 1959--2009: logarithms in per cent, the
unemployment rate in percentage points. Government is real government purchases, income real disposable income. Source: statsmodels macrodata; our computation.}
\end{center}

% figdata: moments (consumption rel 0.80, corr 0.87; investment rel 4.66, corr 0.91; government corr -0.06; unemployment corr -0.88; output ac1 0.86)
Three facts stand out. Consumption moves with output but less, 0.80 times as much; investment moves with it and much more, 4.66
times as much. Unemployment is strongly \omterm{def:b2:business-cycles:procyclical}{countercyclical}, with a correlation of $-0.88$; government purchases are acyclical, at
$-0.06$. And deviations are persistent: output's \omterm{def:b2:business-cycles:cyclical-component}{cyclical component} has a correlation of 0.86 with its value a quarter earlier.

\begin{omchart}{United States: HP cyclical components of real GDP, consumption and investment, quarterly 1959--2009, per
  cent; shaded, recessions.}{statsmodels macrodata; recession dates from the NBER; our computation}
\begin{tikzpicture}
\begin{axis}[omrecord, width=0.72\linewidth, height=0.4\linewidth, xmin=1959, xmax=2009.75, ymin=-30, ymax=20,
    xticklabel style={/pgf/number format/1000 sep=}, ylabel={Per cent from trend},
    legend style={at={(0.5,-0.18)}, anchor=north, legend columns=3, font=\scriptsize}]
  \omshadeperiods{figdata/bachelor-2/28-business-cycles/out/recessions.csv}
  \addplot[omMeth, thin] table[x=t, y=inv] {figdata/bachelor-2/28-business-cycles/out/hp.csv};
  \addlegendentry{investment}
  \addplot[omDef, thick] table[x=t, y=gdp] {figdata/bachelor-2/28-business-cycles/out/hp.csv};
  \addlegendentry{GDP}
  \addplot[omThm, thick, dashed] table[x=t, y=cons] {figdata/bachelor-2/28-business-cycles/out/hp.csv};
  \addlegendentry{consumption}
\end{axis}
\end{tikzpicture}
\end{omchart}

\begin{measure}[How recessions are dated]
% ledger: b2:nber-faq, b2:nber-dates
American recessions are dated by a committee of the National Bureau of Economic Research. It defines a recession as a
significant decline in economic activity that is spread across the economy and lasts more than a few months, and weighs
depth, diffusion and duration together. It looks at monthly series (real personal income less transfers, payroll
employment, real consumption, real manufacturing and trade sales, household employment, industrial production) and at
quarterly real GDP and real gross domestic income. It does not use the rule of two consecutive quarters of falling GDP: the
2001 recession had no such two quarters. It announces turning points only once it is confident of them, between 4 and 21 months
after they occurred. The shaded periods of the chart are its recessions, from peak quarter to trough quarter.
\end{measure}

\section{Trend and cycle}\label{sec:b2:business-cycles:trend}

\begin{definition}[Hodrick--Prescott filter]\label{def:b2:business-cycles:hp-filter}
The \emph{Hodrick--Prescott filter}\index{Hodrick--Prescott filter} splits a series $y_t$ into a trend $\tau_t$ and a \omterm{def:b2:business-cycles:cyclical-component}{cyclical component} $y_t - \tau_t$ by choosing the trend that minimises
$\sum_t(y_t - \tau_t)^2 + \lambda\sum_t\bigl((\tau_{t+1} - \tau_t) - (\tau_t - \tau_{t-1})\bigr)^2$; the larger $\lambda$, the smoother the trend. For quarterly data $\lambda = 1600$ is the usual choice.
\end{definition}

\begin{omfigure}
\begin{tikzpicture}[econ]
\begin{axis}[omrecord, grid=none, xmin=0, xmax=40, ymin=0, ymax=6, width=0.6\linewidth, height=0.36\linewidth,
    xlabel={Time}, ylabel={Log output}, ytick=\empty, legend style={at={(1.03,0.5)}, anchor=west, font=\small}]
  \addplot[omDef, thick, domain=0:40, samples=200] {1 + 0.1*x + 0.5*sin(deg(x/2.2)) + 0.2*sin(deg(x/0.9))};
  \addlegendentry{series}
  \addplot[omThm, thick, dashed, domain=0:40, samples=2] {1 + 0.1*x};
  \addlegendentry{trend}
  \draw[<->, guide] (axis cs:3.45,1.345) -- (axis cs:3.45,1.718);
\end{axis}
\end{tikzpicture}

{\small The decomposition: the \omterm{def:b2:business-cycles:cyclical-component}{cyclical component} is the vertical distance between the series and its trend (arrow). Which
trend is right is not observed; it is chosen by the method.}
\end{omfigure}

The decomposition is a convention, not a measurement. The HP trend bends toward the last observations, so the cycle at the
end of a sample is revised as new data arrive. A filter that removes slow movements can create cycles in a series that has
none: applied to a random walk, it produces a \omterm{def:b2:business-cycles:cyclical-component}{cyclical component} with a regular-looking period. First differences, the
other common choice, keep the slow movements in and amplify quarter-to-quarter noise. Moments of the cycle therefore
depend on the method (Exercise 9).

\section{The real-business-cycle model as a toy model}\label{sec:b2:business-cycles:rbc}

\begin{definition}[Technology shock]\label{def:b2:business-cycles:technology-shock}
A \emph{technology shock}\index{technology shock} is an unexpected change in total factor productivity, the efficiency with
which inputs are turned into output.
\end{definition}

\begin{definition}[Real-business-cycle model]\label{def:b2:business-cycles:rbc}
A \emph{real-business-cycle model}\index{real-business-cycle model} is a stochastic version of the \omterm{def:b2:ramsey-model:ramsey}{Ramsey model} in which
fluctuations arise from \omterm{def:b2:business-cycles:technology-shock}{technology shocks}, markets are competitive and clear, and prices and money play no role.
\end{definition}

The analytic case has log utility, Cobb--Douglas production and full depreciation. A planner chooses consumption and next
period's capital to maximise $E_0\sum_t\beta^t\ln c_t$ subject to $c_t + k_{t+1} = a_tk_t^\alpha$, with $\ln a_t = \rho\ln a_{t-1} + \varepsilon_t$. Guess that a constant share $s$ of output is saved, $k_{t+1} = sy_t$. The \omterm{def:b2:intertemporal-choice:euler-equation}{Euler equation}
$1/c_t = \beta E_t[\alpha y_{t+1}/(k_{t+1}c_{t+1})]$ with $c = (1 - s)y$ becomes $\frac{1}{(1 - s)y_t} = \beta E_t\bigl[\frac{\alpha y_{t+1}}{sy_t(1 - s)y_{t+1}}\bigr] = \frac{\alpha\beta}{s(1 - s)y_t}$, so $s = \alpha\beta$: the guess is verified (Exercise 10).

In per cent deviations from the steady state, $\hat k_{t+1} = \hat y_t$ and $\hat y_t = \hat a_t + \alpha\hat k_t$, so
\begin{equation}\label{eq:b2:business-cycles:propagation}
\hat y_t = \hat a_t + \alpha\hat y_{t-1}.
\end{equation}
Output today is high because technology is high and because capital, built from yesterday's output, is high. The second
term is the model's propagation mechanism.

\begin{definition}[Impulse response]\label{def:b2:business-cycles:impulse-response}
The \emph{impulse response}\index{impulse response} of a variable to a shock is the path of its deviation from what it
would otherwise have been, after a one-time shock of given size and no other shocks.
\end{definition}

The listing computes the \omterm{def:b2:business-cycles:impulse-response}{impulse response} of output to a 1 per cent \omterm{def:b2:business-cycles:technology-shock}{technology shock}, for a shock that disappears at once ($\rho = 0$) and
one that decays at 10 per cent a period ($\rho = 0.9$), and simulates the model's moments.

\omcode{code/bachelor-2/28-business-cycles/listings/rbc.py}{1}{35}{\omterm{def:b2:business-cycles:impulse-response}{Impulse responses} and moments of the analytic
\omterm{def:b2:business-cycles:rbc}{real-business-cycle model}.}\label{lst:b2:business-cycles:rbc}

\omoutput{code/bachelor-2/28-business-cycles/outputs/rbc.txt}{With a one-time shock the response dies in a period or two;
with a persistent shock it has a hump; consumption and investment are exactly as volatile as output.}\label{out:b2:business-cycles:rbc}

\begin{omfigure}
\begin{tikzpicture}[econ]
\begin{axis}[omrecord, grid=none, xmin=-0.3, xmax=12.3, ymin=0, ymax=1.4, width=0.55\linewidth, height=0.36\linewidth,
    xlabel={Periods after the shock}, ylabel={Output, per cent}, legend style={at={(1.03,0.5)}, anchor=west, font=\small}]
  \addplot[omDef, thick, mark=*, mark size=1.3pt, samples at={0,...,12}] {(0.9^(x+1) - (1/3)^(x+1))/(0.9 - 1/3)};
  \addlegendentry{$\rho = 0.9$}
  \addplot[omThm, thick, mark=square*, mark size=1.3pt, samples at={0,...,12}] {(1/3)^x};
  \addlegendentry{$\rho = 0$}
  \addplot[black!50, thin, dashed, samples at={0,...,12}] {0.9^x};
  \addlegendentry{technology, $\rho = 0.9$}
\end{axis}
\end{tikzpicture}

{\small \omterm{def:b2:business-cycles:impulse-response}{Impulse responses} of output to a 1 per cent \omterm{def:b2:business-cycles:technology-shock}{technology shock}, $\alpha = 1/3$: $\hat y_t = (\rho^{t+1} - \alpha^{t+1})/(\rho - \alpha)$. With $\rho = 0$ the response is $\alpha^t$: two thirds gone after one period. With $\rho = 0.9$ output rises above
the shock for a while, as capital accumulates, then follows it down.}
\end{omfigure}

% code: listings/rbc.py
With a one-time shock, output is 0.33 per cent above trend after one period and 0.11 after two: capital accumulation adds
almost no persistence. The persistence of the data must come from the shock itself. With $\rho = 0.9$ output peaks at 1.23 per cent one
period after the shock, and the simulated output has an autocorrelation of 0.95.

\begin{model}{Real-business-cycle model}
\assumes A \omterm{def:b2:ramsey-model:representative-household}{representative household}, competitive markets that always clear, \omterm{def:b2:business-cycles:technology-shock}{technology shocks} as the only disturbance;
in the analytic case log utility, full depreciation and no labour choice.
\predicts \omterm{def:b2:business-cycles:procyclical}{Procyclical} consumption and investment; persistent output if \omterm{def:b2:business-cycles:technology-shock}{technology shocks} are persistent; in versions
with partial depreciation, solved numerically, consumption smoother than output, because households spread a temporary gain
over time, and investment more volatile.
\tested Against the moments of the table. The signs agree, consumption and investment being \omterm{def:b2:business-cycles:procyclical}{procyclical} with correlations
of 0.87 and 0.91 with output, but almost every model of the cycle predicts those signs, so they do not single this one out.
The model's own claims, the relative volatilities and the persistence, are what the table tests, and the analytic case
fails them (next field); the numerical versions, calibrated to match them, are the Year 3 volume's subject.
\breaks The model propagates shocks weakly: persistence must be put into the shocks. In the analytic case consumption and
investment are exactly as volatile as output, against 0.80 and 4.66 in the data. There is no role for money, prices or finance, so the
model has nothing to say about recessions that follow a financial crisis (\cref{ch:b2:banks-and-crises}).
\end{model}

\section{Technology, demand, financial and oil shocks}\label{sec:b2:business-cycles:shocks}

\omterm{def:b2:business-cycles:technology-shock}{Technology shocks} are one candidate. Others are shocks to demand (to spending plans, fiscal policy or monetary policy, which
move output when prices are slow to adjust, as in \cref{ch:b2:three-equation-model}), financial shocks (a fall in banks'
ability or willingness to lend, the subject of the next chapter) and shocks to the price of an input such as oil, which act
like a fall in technology for an oil importer.

\begin{record}[Oil shocks and recessions]
% ledger: b2:hamilton1983
Hamilton (1983) found that all but one of the American recessions after the Second World War had been preceded, typically
by about three quarters of a year, by a dramatic rise in the price of crude oil; the exception was the recession of
1960--61. He stressed that this does not show that oil shocks caused the recessions, but found that over 1948--72 the
correlation was statistically significant and not spurious, which supports the view that oil shocks contributed to at least
some of them.
\end{record}

\section{How a shock is identified}\label{sec:b2:business-cycles:identification}

Output, prices and interest rates move together, so a correlation does not say which shock moved them. Two families of
methods are used. The narrative approach reads the record (central-bank minutes, budget speeches, oil-market events) to
find changes in policy or prices made for reasons unrelated to the current state of the economy, and traces what followed.
Timing restrictions assume that some variables cannot respond within a period to some shocks: for instance, that output
does not respond within the month to a change in the policy rate, so that the month's surprise in the rate can be read as
a policy shock. Both give \omterm{def:b2:business-cycles:impulse-response}{impulse responses}; both rest on assumptions that cannot be fully tested within the data they use.

\section{The fall in volatility after 1985 and its end}\label{sec:b2:business-cycles:moderation}

\begin{definition}[Great Moderation]\label{def:b2:business-cycles:great-moderation}
The \emph{Great Moderation}\index{Great Moderation} is the fall in the volatility of output growth and inflation in the
United States and other rich countries from the mid-1980s to 2007.
\end{definition}

\begin{record}[The Great Moderation and its end]
% ledger: b2:stock-watson2002, b2:mcconnell2000; figdata: moderation (2.47, 1.58, 1.01)
McConnell and Perez-Quiros (2000) dated a break in the volatility of American GDP growth to the first quarter of 1984 and
traced it to durable goods production. Stock and Watson (2002) found that the standard deviation of annual GDP growth fell from 2.7 per cent in 1960--1983 to
1.6 per cent in 1984--2001, and attributed 20 to 30 per cent of the fall to better policy, 20 to 30 per cent to identifiable good luck in productivity
and commodity-price shocks, and 40 to 60 per cent to other good luck. On the data of this chapter the figures are 2.47 and 1.58, and 1.01 for 2002--2008.
The opening chart shows the end: the median volatility of the decade to 2020, 2.10, is back near that of the decade to
1975, 2.21.
\end{record}

\section{Exercises}\label{sec:b2:business-cycles:exercises}

\begin{exercise}[$\star$]\label{exo:b2:business-cycles:1}
Read the opening chart. Give the median volatility for the decades to 1913, 1950, 2007 and 2020. Why is the 1950 value so high?
\end{exercise}

\begin{exercise}[$\star$]\label{exo:b2:business-cycles:2}
From the table: which variables are \omterm{def:b2:business-cycles:procyclical}{procyclical}, which \omterm{def:b2:business-cycles:procyclical}{countercyclical}, which acyclical? Which is most volatile relative to output?
\end{exercise}

\begin{exercise}[$\star$]\label{exo:b2:business-cycles:3}
Compare the standard deviations of annual growth in Stock and Watson's two periods with those computed on this chapter's data.
By what proportion did volatility fall?
\end{exercise}

\begin{exercise}[$\star$]\label{exo:b2:business-cycles:4}
Why does the committee that dates American recessions not use the rule of two quarters of falling GDP? Give the example it
gives.
\end{exercise}

\begin{exercise}[$\star\star$]\label{exo:b2:business-cycles:5}
Predict, then compute: with $\rho = 0$ and $\alpha = 1/3$, how long does it take for the response of output to fall to half its initial size?
\end{exercise}

\begin{exercise}[$\star\star$]\label{exo:b2:business-cycles:6}
With $\alpha = 1/3$ and $\beta = 0.96$, what shares of output are saved and consumed in the analytic model? Why do they not depend on technology?
\end{exercise}

\begin{exercise}[$\star\star$]\label{exo:b2:business-cycles:7}
Modify the listing to set $\rho = 0.5$. Give the response of output in the first four periods, and compare it with the closed form.
\end{exercise}

\begin{exercise}[$\star\star$]\label{exo:b2:business-cycles:8}
Explain why the analytic model gives consumption exactly as volatile as output, and what change to the model makes it
smoother.
\end{exercise}

\begin{exercise}[$\star\star$]\label{exo:b2:business-cycles:9}
\textbf{(D)} With the macrodata, compute the standard deviations of quarterly log growth of consumption and investment relative to
that of output, and their correlations with output growth. Compare with the HP table.
\end{exercise}

\begin{exercise}[$\star\star\star$]\label{exo:b2:business-cycles:10}
Verify that $k_{t+1} = \alpha\beta y_t$ satisfies the \omterm{def:b2:intertemporal-choice:euler-equation}{Euler equation} of the analytic model for every realisation of $a_{t+1}$.
\end{exercise}

\begin{exercise}[$\star\star\star$]\label{exo:b2:business-cycles:11}
From \eqref{eq:b2:business-cycles:propagation} and $\hat a_t = \rho^t$, derive $\hat y_t = (\rho^{t+1} - \alpha^{t+1})/(\rho - \alpha)$, and show that the sum of the responses is $1/((1 - \rho)(1 - \alpha))$.
\end{exercise}

\begin{exercise}[$\star\star\star$]\label{exo:b2:business-cycles:12}
Show that the HP trend equals the series when $\lambda = 0$, and that it tends to a straight line as $\lambda \to \infty$.
\end{exercise}

\section{Problem: Brenland's harvest shocks}\label{sec:b2:business-cycles:problem}

\begin{problem}[{Weekend problem --- Brenland's harvest shocks: the planner's problem, the analytic policy rule, a 1 per
cent shock, and persistence}]
\label{pb:b2:business-cycles:1}
Brenland's output is $y_t = a_tk_t^{0.3}$, where $a_t$ is the harvest; capital lasts one year; the planner maximises $E_0\sum_t 0.95^t\ln c_t$; and $\ln a_t = 0.8\ln a_{t-1} + \varepsilon_t$.

\medskip
\noindent\textbf{Part I --- The planner's problem.}
\begin{enumerate}
  \item Write the resource constraint.
  \item Write the \omterm{def:b2:intertemporal-choice:euler-equation}{Euler equation}.
  \item Why is full depreciation unrealistic for an economy's whole capital stock?
  \item Find the steady-state capital with $a = 1$.
  \item Find steady-state output.
\end{enumerate}

\medskip
\noindent\textbf{Part II --- The analytic policy rule.}
\begin{enumerate}[resume]
  \item Guess $k_{t+1} = sy_t$ and find $s$.
  \item What share of output is consumed?
  \item Write the law of motion of $\hat y_t$.
  \item How does consumption deviate from its steady state?
  \item And investment?
\end{enumerate}

\medskip
\noindent\textbf{Part III --- A 1 per cent shock.} A good harvest raises $a$ by 1 per cent in year 0.
\begin{enumerate}[resume]
  \item What is output in year 0?
  \item In year 1?
  \item In year 2?
  \item When does the response peak, and why does it peak after the shock?
  \item What is the response in year 5?
  \item How does consumption respond?
\end{enumerate}

\medskip
\noindent\textbf{Part IV --- Persistence.}
\begin{enumerate}[resume]
  \item After how many years is the response below half its peak?
  \item With $\rho = 0$, how long would the response take to halve?
  \item What does the comparison say about where persistence comes from?
  \item Find the sum of the responses over all years.
  \item Split it into a factor due to the shock and a factor due to capital.
  \item With $\rho = 0.95$, when does the response peak, and after how many years is it below half the peak?
  \item In the data, output's \omterm{def:b2:business-cycles:cyclical-component}{cyclical component} has an autocorrelation of 0.86 a quarter. What would the model need to match it?
  \item State the result: the \omterm{def:b2:business-cycles:impulse-response}{impulse response} of output and its half-life.
\end{enumerate}
\end{problem}
